% ============================================================================
%  Arithmetic Geometry as the Source Code of Spacetime (English Version)
% ============================================================================
\documentclass[12pt,a4paper]{article}

\usepackage[T1]{fontenc}
\usepackage{amsmath,amssymb,amsthm}
\usepackage{mathrsfs}
\usepackage{graphicx}
\usepackage{hyperref}
\usepackage[margin=2.5cm]{geometry}
\usepackage{booktabs}
\usepackage{enumitem}
\usepackage{xcolor}
\usepackage{tcolorbox}
\usepackage{fancyhdr}
\usepackage{tikz}
\usetikzlibrary{arrows.meta,shapes.geometric,positioning,calc,decorations.pathmorphing}

\tcbuselibrary{skins,breakable}

% Theorem environments
\theoremstyle{definition}
\newtheorem{definition}{Definition}[section]
\newtheorem{conjecture}[definition]{Conjecture}
\newtheorem{proposition}[definition]{Proposition}
\newtheorem{remark}[definition]{Remark}
\newtheorem{analogy}[definition]{Analogy}
\newtheorem{thought}[definition]{Thought Experiment}

% Custom colors
\definecolor{deepblue}{RGB}{26,35,126}
\definecolor{accent}{RGB}{183,28,28}
\definecolor{mathgreen}{RGB}{27,94,32}
\definecolor{physpurple}{RGB}{74,20,140}

% Custom boxes
\newtcolorbox{keyinsight}[1][]{
  colback=deepblue!5, colframe=deepblue!70,
  fonttitle=\bfseries\sffamily, title={#1}, breakable
}
\newtcolorbox{physbox}[1][]{
  colback=physpurple!5, colframe=physpurple!60,
  fonttitle=\bfseries\sffamily, title={#1}, breakable
}
\newtcolorbox{techbox}[1][]{
  colback=accent!5, colframe=accent!60,
  fonttitle=\bfseries\sffamily, title={#1}, breakable
}

\pagestyle{fancy}
\fancyhf{}
\fancyhead[L]{\sffamily\small Arithmetic Geometry as the Source Code of Spacetime}
\fancyhead[R]{\sffamily\small 2026}
\fancyfoot[C]{\thepage}

% Useful commands
\newcommand{\Spec}{\operatorname{Spec}}
\newcommand{\Sh}{\operatorname{Sh}}
\newcommand{\Hom}{\operatorname{Hom}}
\newcommand{\Set}{\mathbf{Set}}
\newcommand{\Topos}{\mathbf{Topos}}
\newcommand{\Zeta}{\mathcal{Z}}

\begin{document}

% ============================================================================
%  TITLE
% ============================================================================
\begin{titlepage}
\centering
\vspace*{2cm}

{\LARGE\sffamily\bfseries
Arithmetic Geometry\\[10pt]
as the Source Code of Spacetime
}

\vspace{0.8cm}

{\Large\sffamily
Towards Quantum Gravity and Spacetime Engineering\\
via Topos Theory, IUT Theory, and Noncommutative Geometry
}

\vspace{1.5cm}

{\large\bfseries Wright Brothers}

\vspace{0.3cm}

{\normalsize\itshape bsbraveshine777@gmail.com}

\vspace{1.5cm}

{\large \today}

\vspace{3cm}

\begin{tcolorbox}[colback=deepblue!3, colframe=deepblue!40, width=13cm]
\sffamily\small
\textbf{Abstract.} Just as Riemannian geometry became the mathematical language of general relativity, modern arithmetic geometry---Grothendieck's topos theory, Mochizuki's Inter-universal Teichm\"{u}ller (IUT) theory, and Connes' noncommutative geometry---may serve as the ``next-generation language'' for quantum gravity and ultimate cosmology. This paper presents a theoretical framework in which arithmetic geometry deciphers the ``source code'' of spacetime, organized along four axes: (1)~topos-theoretic reformulation of quantum mechanics, (2)~cosmological implications of IUT theory, (3)~arithmetic structures in renormalization via noncommutative geometry, and (4)~$p$-adic physics and adelic spacetime. We further develop thought experiments on technological consequences that these theories may enable---direct control of spacetime metrics, topological quantum computation, and vacuum energy extraction.
\end{tcolorbox}

\end{titlepage}

% ============================================================================
%  TABLE OF CONTENTS
% ============================================================================
\tableofcontents
\newpage

% ============================================================================
%  SECTION 1: INTRODUCTION
% ============================================================================
\section{Introduction: Mathematics as the Compiler of Physics}

\subsection{Historical Lesson: The Co-evolution of Geometry and Physical Law}

In the history of the relationship between mathematics and physics, the most dramatic episode is the union of Riemannian geometry with general relativity. The geometry of curved spaces developed by Bernhard Riemann in his 1854 Habilitation lecture ``On the Hypotheses Which Lie at the Foundations of Geometry'' \cite{Riemann1854} was conceived with no physical application in mind. Yet sixty years later, Einstein---with the assistance of Grossmann---adopted Riemannian geometry as the language of gravity and completed general relativity \cite{Einstein1915}.

The lesson of this history is clear: \textbf{the most abstract mathematics can become the descriptive language of the most fundamental physical laws.} And conversely, unsolved problems in physics stimulate the development of new mathematics.

\subsection{The Present Crisis: A Clash of Two Operating Systems}

Modern physics runs on two distinct operating systems (OS): quantum mechanics and general relativity. Quantum mechanics describes the microscopic world; general relativity describes macroscopic spacetime. These two OS ``crash'' in the following situations:

\begin{enumerate}[leftmargin=2em]
  \item \textbf{Black hole singularities}: At the singularities predicted by general relativity (divergence of curvature), quantum effects become dominant, yet no formulation of quantum gravity exists \cite{Penrose1965, Hawking1970}.
  \item \textbf{The moment of cosmic creation}: Describing physics at the Big Bang singularity requires a quantum theory of gravity at the Planck scale ($\ell_P \sim 10^{-35}$\,m, $t_P \sim 10^{-43}$\,s).
  \item \textbf{The information paradox}: The evaporation of black holes via Hawking radiation may violate the unitarity (information conservation) of quantum mechanics \cite{Hawking1976}.
\end{enumerate}

String theory \cite{Polchinski1998} and loop quantum gravity \cite{Rovelli2004} are the leading approaches aiming at this unification, but neither has been experimentally verified. This paper explores a different---or complementary---direction: the possibility that \textbf{arithmetic geometry} provides a ``new language'' for quantum gravity.

\subsection{Central Thesis}

\begin{keyinsight}[Central Thesis]
The theories of modern arithmetic geometry---Grothendieck's topos theory \cite{SGA4}, Mochizuki's IUT theory \cite{Mochizuki2012a}, and Connes' noncommutative geometry \cite{Connes1994}---harbor a rich and still insufficiently explored potential as candidates for a ``meta-language''---a source code of spacetime---that unifies quantum mechanics and general relativity.
\end{keyinsight}

This paper does not present rigorous mathematical proofs. Rather, it systematically catalogues \textbf{structural analogies} between existing mathematical structures and physical problems, and proposes directions for future research programs as a \textbf{thought experiment}.

\subsection{Organization}

Section~\ref{sec:topos} discusses the relationship between topos theory and quantum mechanics. Section~\ref{sec:iut} explores the cosmological implications of IUT theory. Section~\ref{sec:noncommutative} analyzes the arithmetic structures of noncommutative geometry and the renormalization group. Section~\ref{sec:padic} examines $p$-adic physics and adelic spacetime. Section~\ref{sec:technology} develops thought experiments on technological consequences. Section~\ref{sec:conclusion} presents conclusions and outlook.


% ============================================================================
%  SECTION 2: TOPOS THEORY AND QUANTUM MECHANICS
% ============================================================================
\newpage
\section{Reformulating Quantum Mechanics via Topos Theory}
\label{sec:topos}

\subsection{Grothendieck's Topos: A Revolution in the Concept of ``Space''}

The concept of a topos, introduced by Alexandre Grothendieck in the 1960s \cite{SGA4, Grothendieck1957}, demanded a fundamental revision of the answer to the question ``What is a space?''

In classical topology, a space is a ``set of points'' equipped with a family of open subsets (a topology). Grothendieck proposed instead that \textbf{a space should be characterized not as a ``set of points'' but as a ``category of sheaves'' over it}.

\begin{definition}[Grothendieck Topos]
Given a small category $\mathcal{C}$ equipped with a Grothendieck topology $J$, the category of sheaves $\Sh(\mathcal{C}, J)$ on the site $(\mathcal{C}, J)$ is called a \textbf{Grothendieck topos}.
\end{definition}

The revolutionary aspect of this definition is that it allows one to define a ``space'' without presupposing ``points.'' A topos carries its own internal logic, which is not restricted to classical (Boolean) logic but naturally incorporates \textbf{intuitionistic logic} \cite{MacLaneMoerdijk1992, Johnstone2002}.

This property creates the point of contact with quantum mechanics.

\subsection{Quantum Contextuality and Topoi}

One of the most profound features of quantum mechanics is \textbf{contextuality}---the property that the value of a physical quantity depends on the context of measurement. The Kochen--Specker theorem \cite{KochenSpecker1967} proved the impossibility of ``hidden variable'' models that simultaneously assign definite values to all observables of a quantum system.

This demonstrates the limitations of the set-theoretic (Boolean) framework. \textbf{Topos theory provides precisely the mathematical language to transcend this limitation.}

\subsubsection{The Isham--Butterfield Programme}

Chris Isham and Jeremy Butterfield initiated an ambitious programme of reformulating quantum mechanics using topos theory in a series of papers from 1998 to 2002 \cite{IshamButterfield1998, IshamButterfield1999, IshamButterfield2000, IshamButterfield2002}.

Their fundamental insight is as follows:

\begin{keyinsight}[The Isham--Butterfield Insight]
``Propositions'' in quantum mechanics (e.g., ``the energy of the particle lies in the interval $[a,b]$'') should take values not as classical truth values but as elements of the \textbf{subobject classifier} $\Omega$ of a topos. Since $\Omega$ is in general a Heyting algebra rather than a Boolean algebra, it embodies intuitionistic logic in which the law of excluded middle does not hold.
\end{keyinsight}

\subsubsection{D\"{o}ring--Isham's ``What is a Thing?''}

Andreas D\"{o}ring and Isham further developed this programme, publishing a tetralogy entitled ``What is a Thing?'' \cite{DoeringIsham2008a, DoeringIsham2008b, DoeringIsham2008c, DoeringIsham2008d} (a title deliberately echoing Heidegger's work of the same name).

Their construction proceeds as follows:

\begin{enumerate}[leftmargin=2em]
  \item Fix a von Neumann algebra $\mathcal{N}$ (describing the observables of a quantum system).
  \item Consider the poset $\mathcal{V}(\mathcal{N})$ of commutative subalgebras of $\mathcal{N}$ (contexts), ordered by inclusion.
  \item The topos of presheaves $\Set^{\mathcal{V}(\mathcal{N})^{\mathrm{op}}}$ over this poset describes the ``state space'' of the quantum system.
  \item The internal logic of this topos is intuitionistic, naturally reflecting the Kochen--Specker theorem.
\end{enumerate}

The \textbf{spectral presheaf} $\underline{\Sigma}$ in the D\"{o}ring--Isham topos assigns to each context $V \in \mathcal{V}(\mathcal{N})$ the Gelfand spectrum $\Sigma_V$ (a classical state space) of $V$. This presheaf constitutes a ``generalized state space'' for the quantum system.

A proposition ``the value of $\hat{A} \in \mathcal{N}$ lies in $\Delta \subseteq \mathbb{R}$'' is represented as a subobject of the spectral presheaf, and its truth value is given as an element of the Heyting algebra $\Omega$.

\subsubsection{Heunen--Landsman--Spitters' Bohrification}

Chris Heunen, Nicolaas Landsman, and Bas Spitters, independently of D\"{o}ring--Isham, developed the ``Bohrification'' programme inspired by Niels Bohr's complementarity principle \cite{HeunenLandsmanSpitters2009}.

In their approach, given a C*-algebra $A$, one considers the topos over the category of commutative C*-subalgebras and applies Gelfand duality internally within this topos, describing the ``classical aspects'' of the quantum system in the internal logic of the topos.

\begin{proposition}[Fundamental Result of Bohrification \cite{HeunenLandsmanSpitters2009}]
In the Bohrification topos of a C*-algebra $A$, the internal Gelfand spectrum $\underline{\Sigma}_A$ completely recovers the Jordan algebra structure of $A$. That is, the order structure and algebraic operations of the self-adjoint elements of $A$ are reproduced as internal properties of $\underline{\Sigma}_A$ within the topos.
\end{proposition}

\subsection{Topos Quantum Theory and Cosmology: The Problem of the Internal Observer}

The reason the topos-theoretic approach is particularly promising for quantum cosmology is the \textbf{internal observer} problem.

Standard quantum mechanics implicitly assumes a ``classical observer'' external to the system. But in quantum cosmology---the attempt to treat the entire universe as a quantum system---the observer is contained within the system, and no external observer exists \cite{DeWitt1967, Hartle1995}.

Topos theory provides a natural framework for this problem:

\begin{thought}[Topos Quantum Cosmology]
Suppose there exists a topos $\mathcal{E}$ of quantum gravity describing the entire universe. The ``states'' and ``observables'' defined internally in this topos require no external observer. The internal logic of the topos (intuitionistic logic) describes the ``truth'' experienced by observers within the universe. The failure of the law of excluded middle reflects the absence of a ``definite spacetime'' at the quantum gravity scale.
\end{thought}

The \emph{daseinisation} proposed by D\"{o}ring and Isham \cite{DoeringIsham2008a}---named after Heidegger's \emph{Dasein}---is the process of interpreting quantum-mechanical propositions as ``families of approximately valid classical propositions,'' providing a concrete physical interpretation of the internal logic of the topos.

\subsection{Recent Developments}

Topos quantum theory has developed in several directions in recent years:

\begin{itemize}[leftmargin=2em]
  \item \textbf{Flori's systematic formulation}: A comprehensive textbook on topos quantum theory \cite{Flori2013} has been published, advancing the systematization of the field.
  \item \textbf{Sheaf cohomology and quantum contextuality}: Abramsky and Brandenburger \cite{AbramskyBrandenburger2011} demonstrated that quantum non-locality and contextuality admit a unified description in the language of sheaf cohomology, establishing a direct connection between topos-theoretic methods and quantum information theory.
  \item \textbf{Connection to homotopy type theory}: The connection between Univalent Foundations \cite{HoTTbook2013} and topos theory is being explored, suggesting applications of $\infty$-topoi \cite{Lurie2009} to physics. Schreiber \cite{Schreiber2013} formulated gauge theory and Chern--Simons theory within the framework of cohesive $\infty$-topoi, directly implementing Grothendieck's topos-theoretic vision in modern mathematical physics.
  \item \textbf{Factorization algebras in QFT}: Costello \cite{Costello2011} and Costello--Gwilliam \cite{CostelloGwilliam2017} developed the factorization algebra approach, rigorously treating renormalization in the language of derived algebraic geometry and combining locality with homotopy theory.
\end{itemize}


% ============================================================================
%  SECTION 3: IUT THEORY AND COSMOLOGY
% ============================================================================
\newpage
\section{IUT Theory and Cosmology: Communication Between ``Distinct Mathematical Universes''}
\label{sec:iut}

\subsection{Mochizuki's IUT Theory: Overview}

Inter-universal Teichm\"{u}ller (IUT) theory, published by Shinichi Mochizuki in 2012 \cite{Mochizuki2012a, Mochizuki2012b, Mochizuki2012c, Mochizuki2012d}, is a vast mathematical edifice whose primary goal is the proof of the abc conjecture. While the technical details of IUT theory are extraordinarily difficult, the \textbf{ideas} underlying it carry deep implications for physics \cite{Fesenko2015, Fesenko2016}.

The core ideas of IUT theory may be summarized (with great simplification) as follows:

\begin{enumerate}[leftmargin=2em]
  \item \textbf{Setting up multiple ``mathematical universes''}: In standard mathematics, all arguments take place within a single set-theoretic universe. IUT theory sets up \textbf{multiple ``universes'' (Hodge theaters)} in which the additive and multiplicative structures are related in different ways.
  \item \textbf{``Communication'' between universes}: ``Links'' (theta-links, log-links) that transmit information between these distinct universes are constructed. These links \textbf{distort} the relationship between addition and multiplication but preserve certain invariants (e.g., \'etale structures).
  \item \textbf{Control of indeterminacies}: The ``information loss'' (indeterminacy) accompanying inter-universe communication is precisely controlled, ultimately yielding the inequality corresponding to the abc conjecture.
\end{enumerate}

\subsection{Physical Analogies of IUT Theory}

The structure of IUT theory exhibits remarkable similarities with several central problems in theoretical physics.

\subsubsection{The Multiverse and Hodge Theaters}

The landscape problem of string theory \cite{Susskind2003, Bousso2000} shows that string theory admits more than $10^{500}$ distinct vacuum states (and hence ``universes'' with different physical laws).

\begin{analogy}[IUT Theory and the Multiverse]
If one maps IUT theory's Hodge theaters to physical universes and theta-links to inter-universe information channels, IUT theory may be viewed as a \textbf{prototype for an information theory of the multiverse}---communication between universes with different physical laws (mathematical structures).

Just as the relationship between addition and multiplication differs from universe to universe in IUT theory, the fundamental physical constants and symmetry groups may differ across the multiverse. Yet certain ``\'etale'' invariants---geometric and topological properties---may be preserved.
\end{analogy}

\subsubsection{The Holographic Principle and Log-links}

The holographic principle, epitomized by Maldacena's AdS/CFT correspondence \cite{Maldacena1998}, asserts that a $(d+1)$-dimensional gravitational theory is equivalent to a $d$-dimensional conformal field theory on its boundary. The relationship between boundary and bulk in this correspondence has a structural similarity with the ``logarithmic'' transformations in IUT theory's log-links.

\begin{thought}[Log-links and Black Hole Information]
IUT theory's log-links involve operations that transform additive structures into multiplicative structures via the logarithmic function. Physically, the fact that black hole entropy $S = k_B c^3 A / (4G\hbar)$ (the Bekenstein--Hawking formula \cite{Bekenstein1973, Hawking1975}) is proportional to the \textbf{area} of the event horizon itself---not the logarithm of the area---suggests a non-trivial relationship between ``addition'' and ``multiplication'' in gravity.

Whether the indeterminacy control mechanism of IUT theory could serve as a mathematical model for the ``information loss'' in the black hole information paradox is a question worthy of exploration.
\end{thought}

\subsection{\'Etale Fundamental Groups and Spacetime Topology}

The \'etale fundamental group (arithmetic fundamental group) of algebraic curves plays a central role in IUT theory. The \'etale fundamental group generalizes the theory of covering spaces in algebraic geometry and is deeply connected to Galois groups \cite{SGA1}.

In physics, the topology of spacetime (including the fundamental group) describes non-trivial aspects of gravity. Wormholes, cosmic strings, and topological defects are all related to non-trivial fundamental groups of spacetime \cite{VisserBook1995}.

\begin{conjecture}[Arithmetic Fundamental Group of Spacetime]
The ``fundamental group'' of spacetime at the Planck scale should be described not as an ordinary topological fundamental group but as an \textbf{arithmetic fundamental group} (a structure analogous to the \'etale fundamental group) containing number-theoretic information. In this case, the microstructure of spacetime is described not as a continuous manifold but as an algebro-geometric scheme (or stack).
\end{conjecture}

\subsection{Anabelian Geometry and the Recovery of Physical Laws}

Anabelian geometry, proposed by Grothendieck in his ``Esquisse d'un Programme'' \cite{Grothendieck1997} and substantially developed by Mochizuki \cite{Mochizuki1996}, asserts that \textbf{algebraic curves can be essentially recovered from their \'etale fundamental groups}. That is, group-theoretic and category-theoretic data completely determine geometric objects.

\begin{analogy}[Anabelian Geometry and the Recovery of Physical Laws]
Just as geometry is recovered from groups in anabelian geometry, might physical laws be completely recoverable from the ``arithmetic symmetry group'' of spacetime? That is, the correct formulation of quantum gravity may take as its starting point a certain ``fundamental group'' of spacetime and derive from it the geometry of spacetime and the dynamics of matter fields.
\end{analogy}


% ============================================================================
%  SECTION 4: NONCOMMUTATIVE GEOMETRY
% ============================================================================
\newpage
\section{Noncommutative Geometry and Renormalization: Quantum Fields as Arithmetic Structures}
\label{sec:noncommutative}

\subsection{Connes' Noncommutative Geometry}

Noncommutative geometry, founded by Alain Connes \cite{Connes1994}, is a mathematical framework that extends the concept of ``space'' to noncommutative algebras. An ordinary topological space $X$ is completely recovered by its commutative C*-algebra $C(X)$ of continuous functions (Gelfand--Naimark theorem). Connes' noncommutative geometry extends this correspondence to noncommutative C*-algebras, developing a geometry of ``noncommutative spaces.''

The applications of noncommutative geometry to physics are manifold; this paper focuses on the following three points of contact.

\subsection{The Connes--Kreimer Hopf-Algebraic Renormalization}

Renormalization in quantum field theory---the systematic procedure for handling divergences (infinities) arising in calculations---was long regarded as a mathematically unsatisfactory ``art.'' Connes and Dirk Kreimer discovered that the mathematical structure of renormalization is beautifully described within the framework of Hopf algebras \cite{ConnesKreimer1998, ConnesKreimer2000, ConnesKreimer2001}.

\begin{keyinsight}[The Connes--Kreimer Result]
The Hopf algebra $\mathcal{H}$ defined on the set of Feynman diagrams, together with its Birkhoff decomposition, algebraically describes the action of the renormalization group. This structure has a deep connection with the \textbf{motivic Galois group} in algebraic geometry.
\end{keyinsight}

Specifically, the Birkhoff decomposition of renormalization is equivalent to a Riemann--Hilbert problem for loop groups on the projective line $\mathbb{P}^1(\mathbb{C})$, a central object of algebraic geometry \cite{ConnesMarcolli2008}.

\subsection{The Connes--Marcolli Arithmetic Quantum Mechanics}

In their magnum opus ``Noncommutative Geometry, Quantum Fields and Motives'' \cite{ConnesMarcolli2008}, Connes and Matilde Marcolli systematically developed the deep connections between number theory and quantum field theory.

Of particular note is the \textbf{Bost--Connes system} \cite{BostConnes1995}, a remarkable construction in which special values of the Riemann zeta function emerge naturally as phase transitions of a quantum statistical mechanical system:

\begin{enumerate}[leftmargin=2em]
  \item A C*-dynamical system $(\mathcal{A}, \sigma_t)$ is constructed, where $\mathcal{A}$ is a noncommutative algebra related to the Heisenberg group.
  \item This system has a phase transition at inverse temperature $\beta = 1$.
  \item Evaluating the zeta function using KMS states (thermal equilibrium states) in the low-temperature phase $\beta > 1$ yields the values $\zeta(n)$ ($n \in \mathbb{N}$).
  \item The symmetry group of this system is $\hat{\mathbb{Z}}^* \cong \mathrm{Gal}(\mathbb{Q}^{\mathrm{ab}}/\mathbb{Q})$, related to the reciprocity law of class field theory.
\end{enumerate}

\begin{physbox}[Physical Implications]
The Bost--Connes system shows that fundamental objects of number theory (zeta functions, Galois groups) are described in the natural language of quantum statistical mechanics. This suggests the possibility that \textbf{the distribution of primes and the microstructure of spacetime may be described within the same mathematical framework}.
\end{physbox}

\subsection{The Renormalization Group and Motives: Resolution Transformations of the Universe}

Connes and Marcolli explored the possibility of interpreting the renormalization group as the action of a ``motivic Galois group'' \cite{ConnesMarcolli2008, ConnesMarcolli2004}. Indeed, F.~Brown \cite{Brown2012} proved that the periods of mixed Tate motives over $\Spec(\mathbb{Z})$ are exactly the multiple zeta values (MZVs), establishing that the values appearing in perturbative quantum field theory calculations are fundamental objects of arithmetic geometry. Furthermore, Brown \cite{Brown2017} made the Connes--Marcolli ``cosmic Galois group'' precise for a broad class of Feynman amplitudes, showing they obey a coaction principle governed by motivic structures. Bloch--Esnault--Kreimer \cite{BlochEsnaultKreimer2006} and Marcolli \cite{Marcolli2010} have also deeply studied the structure of motives associated with graph polynomials.

Motives were envisioned by Grothendieck as a ``universal'' structure underlying the various cohomology theories of algebraic geometry (de Rham, \'etale, crystalline, etc.) \cite{Grothendieck1969}. The theory of motives remains incomplete, but Voevodsky's category of mixed motives \cite{Voevodsky2000} is a central object of algebraic geometry.

\begin{analogy}[The Renormalization Group and the Motivic Galois Group]
The action of the renormalization group in quantum field theory---the symmetry that the essence of physics is invariant under changes of energy scale (resolution)---corresponds to the action of the motivic Galois group in algebraic geometry.

Rephrased physically: \textbf{the transformation rules between descriptions of the universe at different scales are described as Galois symmetries of arithmetic geometry.} The renormalization group flow is a ``compilation'' process from the Planck scale (UV) to the macroscopic scale (IR), and the ``compiler'' controlling this process is the motivic Galois group.
\end{analogy}

\subsection{Deninger's Programme: Zeta Functions and Dynamical Systems}

Christopher Deninger has pursued a sweeping programme of interpreting the zeros of arithmetic zeta functions as periodic orbits of dynamical systems \cite{Deninger1998, Deninger2002}.

The Riemann hypothesis asserts that all non-trivial zeros of $\zeta(s)$ lie on the line $\mathrm{Re}(s) = 1/2$. Deninger seeks a framework for understanding this conjecture as ``the distribution of periodic orbits of an appropriate dynamical system on $\Spec(\mathbb{Z})$.''

\begin{thought}[The Riemann Hypothesis and Spacetime Stability]
If Deninger's programme is completed and the Riemann hypothesis is proved in the language of dynamical systems, then the regularity of prime distribution (the Riemann hypothesis) could be interpreted as a \textbf{stability condition for the microstructure of spacetime}. The failure of the Riemann hypothesis would signify the existence of ``unstable modes'' in spacetime at the Planck scale.
\end{thought}


% ============================================================================
%  SECTION 5: P-ADIC PHYSICS
% ============================================================================
\newpage
\section{$p$-Adic Physics and Adelic Spacetime}
\label{sec:padic}

\subsection{$p$-Adic Numbers and Physics}

The $p$-adic number field $\mathbb{Q}_p$ is obtained as an alternative completion of the rationals (using a different absolute value from the standard one that gives $\mathbb{R}$). By Ostrowski's theorem, the non-trivial absolute values on $\mathbb{Q}$ are exhausted by the standard absolute value and the $p$-adic absolute values (one for each prime $p$).

In 1987, I.~V.~Volovich proposed in ``Number Theory as the Ultimate Physical Theory'' \cite{Volovich1987} that spacetime at sub-Planckian scales might have a $p$-adic structure. His basic claim is:

\begin{keyinsight}[The Volovich Hypothesis \cite{Volovich1987}]
Below the Planck length $\ell_P$, the coordinates of spacetime should be described not as elements of $\mathbb{R}$ but as elements of $\mathbb{Q}_p$. The Planck length corresponds to the ``boundary'' between real and $p$-adic descriptions.
\end{keyinsight}

\subsection{$p$-Adic String Theory}

Following Volovich's proposal, P.~G.~O.~Freund and E.~Witten \cite{FreundWitten1987} and Freund and M.~Olson \cite{FreundOlson1987} developed $p$-adic string theory.

The $p$-adic Virasoro amplitude is:
\begin{equation}
  A_p(s,t) = \frac{\zeta_p(1-s)\,\zeta_p(1-t)\,\zeta_p(1-u)}{\zeta_p(s)\,\zeta_p(t)\,\zeta_p(u)}
\end{equation}
where $\zeta_p(s) = (1 - p^{-s})^{-1}$ is the $p$-adic zeta factor and $s + t + u = 0$ is the Mandelstam relation.

Remarkably, the product of the standard string's Veneziano amplitude $A_\infty(s,t)$ and all $p$-adic amplitudes satisfies the \textbf{adelic product formula} \cite{FreundWitten1987}:
\begin{equation}
  A_\infty(s,t) \cdot \prod_{p\,:\,\text{prime}} A_p(s,t) = 1
\end{equation}

This formula demonstrates that string scattering amplitudes possess an \textbf{adelic} structure spanning the real field and all $p$-adic fields.

\subsection{Adelic Spacetime}

The adele ring $\mathbb{A}_\mathbb{Q} = \mathbb{R} \times \prod'_p \mathbb{Q}_p$ (where $\prod'$ denotes the restricted product) is the object simultaneously containing all completions of the rationals.

Yu.~I.~Manin \cite{Manin1989} and B.~Dragovich et al.~\cite{Dragovich2009, Dragovich2017} have explored the possibility that spacetime possesses an adelic structure. Dragovich \cite{Dragovich1995} also succeeded in constructing the adelic harmonic oscillator, showing its spectrum to be related to the zeros of the Riemann zeta function. Manin and Marcolli \cite{ManinMarcolli2014} developed attempts to apply arithmetic-geometric concepts such as modular curves and dessins d'enfants directly to cosmological models:

\begin{thought}[The Adelic Spacetime Hypothesis]
Physical spacetime should be described not as $\mathbb{R}^{3,1}$ but as an adelic space such as $\mathbb{A}_\mathbb{Q}^{3,1}$. The real spacetime $\mathbb{R}^{3,1}$ that we experience macroscopically is merely the ``Archimedean component'' of adelic spacetime; at the Planck scale, the $p$-adic components become manifest.

In this picture, the $p$-adic component corresponding to each prime $p$ constitutes a distinct ``channel'' of spacetime's microstructure. The infinitude of primes corresponds to the infinite richness of spacetime's microstructure.
\end{thought}

\subsection{Spacetime as $\Spec(\mathbb{Z})$}

In arithmetic geometry, the spectrum of the ring of integers $\Spec(\mathbb{Z})$ is the most fundamental object. The ``points'' of $\Spec(\mathbb{Z})$ consist of primes $p$ and the ``generic point'' $\eta$ (corresponding to the field of rationals $\mathbb{Q}$), each carrying a local ring.

Manin \cite{Manin1995} and Connes \cite{Connes1999} have suggested that the geometry of $\Spec(\mathbb{Z})$ may be deeply related to the fundamental framework of physics.

\begin{conjecture}[The $\Spec(\mathbb{Z})$ Spacetime Hypothesis]
The microstructure of spacetime at the Planck scale possesses an arithmetic structure analogous to the scheme $\Spec(\mathbb{Z})$ (or some extension thereof). The usual differentiable-geometric spacetime $(\mathcal{M}, g_{\mu\nu})$ is merely the ``real realization'' of this arithmetic structure (the fiber over $\Spec(\mathbb{Z}) \to \Spec(\mathbb{R})$).

Primes $p$ correspond to the ``fundamental vibrational modes'' of spacetime's microstructure, and the geometry of $\Spec(\mathbb{Z})$ (\'etale cohomology, zeta functions) determines the quantum-gravitational properties of spacetime.
\end{conjecture}


% ============================================================================
%  SECTION 6: TECHNOLOGY
% ============================================================================
\newpage
\section{Technological Consequences: Spacetime Engineering}
\label{sec:technology}

If the theoretical frameworks above are completed and the ``source code'' of spacetime is deciphered, what technologies might become possible? In this section we develop thought experiments that prioritize imagination over rigor.

\subsection{Direct Control of the Spacetime Metric: Non-Newtonian Propulsion}

\subsubsection{An Arithmetic-Geometric Reinterpretation of the Alcubierre Drive}

The ``warp drive'' proposed by Alcubierre in 1994 \cite{Alcubierre1994} achieves superluminal travel by contracting spacetime ahead of the craft and expanding it behind, leaving the craft locally at rest:
\begin{equation}
  ds^2 = -dt^2 + (dx - v_s(t)\, f(r_s)\, dt)^2 + dy^2 + dz^2
\end{equation}
where $v_s(t)$ is the velocity of the warp bubble and $f(r_s)$ is the shape function.

This metric is an exact solution of general relativity but requires negative energy density (exotic matter), making realization difficult \cite{Pfenning1997}. However, new solutions relaxing energy conditions have recently been found by Lentz \cite{Lentz2021} and Bobrick--Martire \cite{BobrickMartire2021}.

\begin{techbox}[Arithmetic-Geometric Warp Drive]
If spacetime possesses $\Spec(\mathbb{Z})$-like arithmetic structure, then spacetime ``distortion'' is described not by the usual Riemann curvature tensor but by arithmetic invariants (e.g., the zero distribution of local zeta functions). In this case, the ``negative energy'' required for a warp drive might be achieved through local modulation of spacetime's arithmetic structure---for example, shifting the zeros of a local zeta function outside the critical strip.

This represents a paradigm of controlling spacetime through switching of discrete arithmetic parameters (informational operations) rather than controlling a continuous metric (requiring enormous energy), potentially enabling dramatic reduction in the energy requirements.
\end{techbox}

\subsection{The Ultimate Form of Topological Quantum Computing}

\subsubsection{Current Topological Quantum Computation}

Topological quantum computation, proposed by Kitaev \cite{Kitaev2003}, uses braiding operations on non-abelian anyons for quantum computation. The topological phase factors accompanying anyon exchange are stable against local perturbations, enabling quantum computation extremely robust to noise.

Freedman, Kitaev, Larsen, and Wang \cite{FreemanKitaevLarsenWang2003} proved that this approach achieves universal quantum computation (BQP-completeness).

\subsubsection{Topos-Theoretic Quantum Computing}

\begin{thought}[Topos Quantum Computation]
While current topological quantum computation exploits topological properties of two-dimensional matter (braid groups of anyons), topos-theoretic quantum computation would directly exploit \textbf{higher categorical structures}.

Specifically, ``higher anyons'' defined within $(\infty,n)$-categories \cite{Lurie2009} ($\infty$-topoi)---algebraic structures corresponding to $n$-dimensional topological field theories under the cobordism hypothesis \cite{BaezDolan1995, Lurie2009Hopkins}---would serve as computational resources.

In this approach, computational error tolerance is guaranteed not by topological invariance but by \textbf{categorical universality} (existence of limits, functoriality, coherence of natural transformations). This represents a natural generalization of current topological quantum computation, potentially possessing vastly greater expressive power.
\end{thought}

\subsection{Vacuum Energy Extraction and Phase Transitions of Spacetime}

\subsubsection{The Casimir Effect and Vacuum Energy}

The Casimir effect \cite{Casimir1948}---the force arising because the vacuum energy density between two conducting plates differs from the exterior---is experimental evidence that the quantum vacuum is an energy-bearing entity \cite{Lamoreaux1997}. The dynamical Casimir effect \cite{Wilson2011} experimentally confirmed that rapidly moving boundaries generate photons from the vacuum.

\subsubsection{Arithmetic Vacuum Structure}

\begin{thought}[Number-Theoretic Vacuum Energy]
If spacetime's microstructure is $\Spec(\mathbb{Z})$-like, then the energy spectrum of the quantum vacuum is related to the distribution of primes. If the zeros of the Riemann zeta function correspond to energy levels of vacuum modes, then the analytic properties of the zeta function (functional equation, zero distribution) determine the structure of vacuum energy.

In this picture, energy extraction from the vacuum could be achieved as a ``phase transition'' of spacetime's arithmetic structure---for example, modulation of $p$-adic structure inducing local changes in the zeta function. The phase transition mechanism of the Bost--Connes system \cite{BostConnes1995} could serve as a prototype for this process.
\end{thought}


% ============================================================================
%  SECTION 7: CONCLUSION
% ============================================================================
\newpage
\section{Conclusions and Outlook}
\label{sec:conclusion}

\subsection{Summary}

This paper has argued from four axes that arithmetic geometry may provide a theoretical framework for deciphering the ``source code'' of spacetime:

\begin{enumerate}[leftmargin=2em]
  \item \textbf{Topos theory} (\S\ref{sec:topos}): The prior work of Isham--Butterfield, D\"{o}ring--Isham, and Heunen--Landsman--Spitters has shown that the internal logic of topoi naturally describes the contextuality of quantum mechanics. Topos quantum theory is one of the most promising language candidates for quantum cosmology that does not presuppose an external observer.

  \item \textbf{IUT theory} (\S\ref{sec:iut}): The structure of ``communication between multiple mathematical universes'' in Mochizuki's IUT theory has deep structural similarities with multiverse physics and the holographic principle. The indeterminacy control mechanism of IUT theory may serve as a mathematical model for the black hole information paradox.

  \item \textbf{Noncommutative geometry} (\S\ref{sec:noncommutative}): The Connes--Kreimer renormalization theory and Connes--Marcolli arithmetic quantum mechanics concretely demonstrate that the structure of quantum field theory is intimately related to the deep structures of arithmetic geometry (motives, Galois groups, zeta functions).

  \item \textbf{$p$-Adic physics} (\S\ref{sec:padic}): The pioneering work of Volovich, Freund--Witten, and Manin has indicated the possibility that Planck-scale spacetime possesses adelic structure, paving the way toward the grand hypothesis that models the microstructure of spacetime on $\Spec(\mathbb{Z})$.
\end{enumerate}

\subsection{Prospect of a Paradigm Shift}

The most ambitious claim of this paper is the possibility of the following paradigm shift:

\begin{keyinsight}[Paradigm Shift]
\textbf{From ``how to move matter and energy'' to ``how to rewrite spacetime itself.''}

If arithmetic geometry becomes the tool for deciphering the ``source code'' of spacetime, physics will transition from the stage of ``manipulating matter on the stage of spacetime'' to the stage of ``designing and modifying the stage itself.''
\end{keyinsight}

\subsection{Future Research Directions}

To develop the structural analogies presented in this paper into rigorous mathematical and physical results, the following research directions are important:

\begin{enumerate}[leftmargin=2em]
  \item \textbf{Construction of concrete models of topos quantum gravity}: Formulate spin foam models and spin networks of loop quantum gravity \cite{Rovelli2004} within the internal logic of a topos.
  \item \textbf{Physical formulation of IUT theory}: Give concrete physical interpretations to the mathematical structures of IUT theory (Hodge theaters, theta-links, log-links).
  \item \textbf{Quantum field theory on $\Spec(\mathbb{Z})$}: Formulate a ``quantum field theory'' on $\Spec(\mathbb{Z})$ and show that its low-energy limit reproduces ordinary quantum field theory.
  \item \textbf{Cosmological applications of the Bost--Connes system}: Relate the phase transitions of the Bost--Connes system to cosmological phase transitions (inflation, spontaneous symmetry breaking).
  \item \textbf{Development of adelic string theory}: Extend the Freund--Witten adelic product formula to the modern formulation of string theory (M-theory).
\end{enumerate}

\subsection{Closing Remarks}

Arithmetic geometry appears, at first glance, to be the product of pure thought most distant from physical reality. Yet, as Riemannian geometry demonstrated, the most abstract mathematics can become the language of the most fundamental physical laws. The mathematical universes opened by Grothendieck's topos theory, Mochizuki's IUT theory, and Connes' noncommutative geometry present, to the cutting edge of humanity's intellectual adventure---quantum gravity and spacetime engineering---a path not yet sufficiently explored.

In the spirit of the project name ``Wright Brothers,'' this path---like the aviation that began with wind tunnel experiments in a bicycle shop---may begin with pencil-and-paper calculations in pure mathematics and eventually lead to technology that controls spacetime itself. The first step is to construct a ``translation dictionary'' between arithmetic geometry and physics.


% ============================================================================
%  REFERENCES
% ============================================================================
\newpage
\begin{thebibliography}{99}

% --- Foundations ---
\bibitem{Riemann1854}
B.~Riemann, ``\"{U}ber die Hypothesen, welche der Geometrie zu Grunde liegen,'' Habilitationsschrift, G\"{o}ttingen (1854).

\bibitem{Einstein1915}
A.~Einstein, ``Die Feldgleichungen der Gravitation,'' \textit{Sitzungsber. Preuss. Akad. Wiss.} (1915), pp.~844--847.

% --- Quantum gravity and black holes ---
\bibitem{Penrose1965}
R.~Penrose, ``Gravitational collapse and space-time singularities,'' \textit{Phys. Rev. Lett.} \textbf{14}, 57--59 (1965).

\bibitem{Hawking1970}
S.~W.~Hawking and R.~Penrose, ``The singularities of gravitational collapse and cosmology,'' \textit{Proc. R. Soc. Lond. A} \textbf{314}, 529--548 (1970).

\bibitem{Hawking1975}
S.~W.~Hawking, ``Particle creation by black holes,'' \textit{Commun. Math. Phys.} \textbf{43}, 199--220 (1975).

\bibitem{Hawking1976}
S.~W.~Hawking, ``Breakdown of predictability in gravitational collapse,'' \textit{Phys. Rev. D} \textbf{14}, 2460--2473 (1976).

\bibitem{Bekenstein1973}
J.~D.~Bekenstein, ``Black holes and entropy,'' \textit{Phys. Rev. D} \textbf{7}, 2333--2346 (1973).

\bibitem{DeWitt1967}
B.~S.~DeWitt, ``Quantum theory of gravity. I. The canonical theory,'' \textit{Phys. Rev.} \textbf{160}, 1113--1148 (1967).

\bibitem{Hartle1995}
J.~B.~Hartle, ``Spacetime quantum mechanics and the quantum mechanics of spacetime,'' in \textit{Gravitation and Quantizations} (Les Houches 1992), North-Holland (1995). arXiv:gr-qc/9304006.

% --- String theory and holography ---
\bibitem{Polchinski1998}
J.~Polchinski, \textit{String Theory}, Vols.~1 \& 2, Cambridge Univ.~Press (1998).

\bibitem{Rovelli2004}
C.~Rovelli, \textit{Quantum Gravity}, Cambridge Univ.~Press (2004).

\bibitem{Maldacena1998}
J.~Maldacena, ``The large $N$ limit of superconformal field theories and supergravity,'' \textit{Adv. Theor. Math. Phys.} \textbf{2}, 231--252 (1998). arXiv:hep-th/9711200.

\bibitem{Susskind2003}
L.~Susskind, ``The anthropic landscape of string theory,'' arXiv:hep-th/0302219 (2003).

\bibitem{Bousso2000}
R.~Bousso and J.~Polchinski, ``Quantization of four-form fluxes and dynamical neutralization of the cosmological constant,'' \textit{JHEP} \textbf{06}, 006 (2000).

% --- Grothendieck and topos theory ---
\bibitem{SGA4}
M.~Artin, A.~Grothendieck, and J.-L.~Verdier, \textit{Th\'{e}orie des Topos et Cohomologie \'{E}tale des Sch\'{e}mas} (SGA 4), Lecture Notes in Math.~269, 270, 305, Springer (1972--1973).

\bibitem{SGA1}
A.~Grothendieck, \textit{Rev\^{e}tements \'{E}tales et Groupe Fondamental} (SGA 1), Lecture Notes in Math.~224, Springer (1971).

\bibitem{Grothendieck1957}
A.~Grothendieck, ``Sur quelques points d'alg\`{e}bre homologique,'' \textit{T\^{o}hoku Math. J.} \textbf{9}, 119--221 (1957).

\bibitem{Grothendieck1969}
A.~Grothendieck, ``Standard conjectures on algebraic cycles,'' in \textit{Algebraic Geometry} (Bombay 1968), Oxford Univ.~Press (1969), pp.~193--199.

\bibitem{Grothendieck1997}
A.~Grothendieck, ``Esquisse d'un programme'' (1984), in \textit{Geometric Galois Actions}, Vol.~1, Cambridge Univ.~Press (1997), pp.~5--48.

\bibitem{MacLaneMoerdijk1992}
S.~Mac Lane and I.~Moerdijk, \textit{Sheaves in Geometry and Logic}, Springer (1992).

\bibitem{Johnstone2002}
P.~T.~Johnstone, \textit{Sketches of an Elephant: A Topos Theory Compendium}, Vols.~1 \& 2, Oxford Univ.~Press (2002).

% --- Topos quantum theory ---
\bibitem{IshamButterfield1998}
C.~J.~Isham and J.~Butterfield, ``A topos perspective on the Kochen--Specker theorem: I.,'' \textit{Int. J. Theor. Phys.} \textbf{37}, 2669--2733 (1998). arXiv:quant-ph/9803055.

\bibitem{IshamButterfield1999}
J.~Butterfield and C.~J.~Isham, ``A topos perspective on the Kochen--Specker theorem: II.,'' \textit{Int. J. Theor. Phys.} \textbf{38}, 827--859 (1999). arXiv:quant-ph/9808067.

\bibitem{IshamButterfield2000}
J.~Butterfield and C.~J.~Isham, ``Some possible roles for topos theory in quantum theory and quantum gravity,'' \textit{Found. Phys.} \textbf{30}, 1707--1735 (2000). arXiv:gr-qc/9910005.

\bibitem{IshamButterfield2002}
J.~Butterfield and C.~J.~Isham, ``A topos perspective on the Kochen--Specker theorem: IV.,'' \textit{Int. J. Theor. Phys.} \textbf{41}, 613--639 (2002). arXiv:quant-ph/0107030.

\bibitem{DoeringIsham2008a}
A.~D\"{o}ring and C.~J.~Isham, ``A topos foundation for theories of physics: I.,'' \textit{J. Math. Phys.} \textbf{49}, 053515 (2008). arXiv:quant-ph/0703060.

\bibitem{DoeringIsham2008b}
A.~D\"{o}ring and C.~J.~Isham, ``A topos foundation for theories of physics: II.,'' \textit{J. Math. Phys.} \textbf{49}, 053516 (2008). arXiv:quant-ph/0703062.

\bibitem{DoeringIsham2008c}
A.~D\"{o}ring and C.~J.~Isham, ``A topos foundation for theories of physics: III.,'' \textit{J. Math. Phys.} \textbf{49}, 053517 (2008). arXiv:quant-ph/0703064.

\bibitem{DoeringIsham2008d}
A.~D\"{o}ring and C.~J.~Isham, ``A topos foundation for theories of physics: IV.,'' \textit{J. Math. Phys.} \textbf{49}, 053518 (2008). arXiv:quant-ph/0703066.

\bibitem{HeunenLandsmanSpitters2009}
C.~Heunen, N.~P.~Landsman, and B.~Spitters, ``A topos for algebraic quantum theory,'' \textit{Commun. Math. Phys.} \textbf{291}, 63--110 (2009). arXiv:0709.4364.

\bibitem{KochenSpecker1967}
S.~Kochen and E.~P.~Specker, ``The problem of hidden variables in quantum mechanics,'' \textit{J. Math. Mech.} \textbf{17}, 59--87 (1967).

\bibitem{Flori2013}
C.~Flori, \textit{A First Course in Topos Quantum Theory}, Lecture Notes in Phys.~868, Springer (2013).

% --- Homotopy type theory and higher categories ---
\bibitem{HoTTbook2013}
The Univalent Foundations Program, \textit{Homotopy Type Theory}, Institute for Advanced Study (2013).

\bibitem{Lurie2009}
J.~Lurie, \textit{Higher Topos Theory}, Annals of Math.~Studies 170, Princeton Univ.~Press (2009).

\bibitem{Lurie2009Hopkins}
J.~Lurie, ``On the classification of topological field theories,'' \textit{Current Dev. Math.} (2008), pp.~129--280. arXiv:0905.0465.

\bibitem{BaezDolan1995}
J.~C.~Baez and J.~Dolan, ``Higher-dimensional algebra and topological quantum field theory,'' \textit{J. Math. Phys.} \textbf{36}, 6073--6105 (1995). arXiv:q-alg/9503002.

\bibitem{CostelloGwilliam2017}
K.~Costello and O.~Gwilliam, \textit{Factorization Algebras in Quantum Field Theory}, Vols.~1 \& 2, Cambridge Univ.~Press (2017, 2021).

% --- IUT Theory ---
\bibitem{Mochizuki2012a}
S.~Mochizuki, ``Inter-universal Teichm\"{u}ller theory I: Construction of Hodge theaters,'' RIMS Preprint 1756 (2012); \textit{PRIMS} \textbf{57}, 3--207 (2021).

\bibitem{Mochizuki2012b}
S.~Mochizuki, ``Inter-universal Teichm\"{u}ller theory II: Hodge--Arakelov-theoretic evaluation,'' RIMS Preprint 1757 (2012); \textit{PRIMS} \textbf{57}, 209--401 (2021).

\bibitem{Mochizuki2012c}
S.~Mochizuki, ``Inter-universal Teichm\"{u}ller theory III: Canonical splittings of the log-theta-lattice,'' RIMS Preprint 1758 (2012); \textit{PRIMS} \textbf{57}, 403--626 (2021).

\bibitem{Mochizuki2012d}
S.~Mochizuki, ``Inter-universal Teichm\"{u}ller theory IV: Log-volume computations and set-theoretic foundations,'' RIMS Preprint 1759 (2012); \textit{PRIMS} \textbf{57}, 627--723 (2021).

\bibitem{Mochizuki1996}
S.~Mochizuki, ``The profinite Grothendieck conjecture for closed hyperbolic curves over number fields,'' \textit{J. Math. Sci. Univ. Tokyo} \textbf{3}, 571--627 (1996).

\bibitem{Fesenko2015}
I.~Fesenko, ``Arithmetic deformation theory via arithmetic fundamental groups and nonarchimedean theta-functions,'' \textit{Eur. J. Math.} \textbf{1}, 405--440 (2015).

\bibitem{Fesenko2016}
I.~Fesenko, ``Fukugen,'' IARA preprint (2016).

% --- Noncommutative geometry ---
\bibitem{Connes1994}
A.~Connes, \textit{Noncommutative Geometry}, Academic Press (1994).

\bibitem{Connes1999}
A.~Connes, ``Trace formula in noncommutative geometry and the zeros of the Riemann zeta function,'' \textit{Selecta Math.} \textbf{5}, 29--106 (1999). arXiv:math/9811068.

\bibitem{ConnesKreimer1998}
A.~Connes and D.~Kreimer, ``Hopf algebras, renormalization and noncommutative geometry,'' \textit{Commun. Math. Phys.} \textbf{199}, 203--242 (1998). arXiv:hep-th/9808042.

\bibitem{ConnesKreimer2000}
A.~Connes and D.~Kreimer, ``Renormalization in QFT and the Riemann--Hilbert problem I,'' \textit{Commun. Math. Phys.} \textbf{210}, 249--273 (2000). arXiv:hep-th/9912092.

\bibitem{ConnesKreimer2001}
A.~Connes and D.~Kreimer, ``Renormalization in QFT and the Riemann--Hilbert problem II,'' \textit{Commun. Math. Phys.} \textbf{216}, 215--241 (2001). arXiv:hep-th/0003188.

\bibitem{ConnesMarcolli2008}
A.~Connes and M.~Marcolli, \textit{Noncommutative Geometry, Quantum Fields and Motives}, AMS Colloquium Publications (2008).

\bibitem{BostConnes1995}
J.-B.~Bost and A.~Connes, ``Hecke algebras, type III factors and phase transitions with spontaneous symmetry breaking in number theory,'' \textit{Selecta Math.} \textbf{1}, 411--457 (1995).

% --- Motives ---
\bibitem{Voevodsky2000}
V.~Voevodsky, ``Triangulated categories of motives over a field,'' Annals of Math.~Studies 143, Princeton Univ.~Press (2000).

% --- Deninger ---
\bibitem{Deninger1998}
C.~Deninger, ``Some analogies between number theory and dynamical systems on foliated spaces,'' in \textit{Proc. ICM Berlin 1998}, Vol.~I, pp.~163--186.

\bibitem{Deninger2002}
C.~Deninger, ``A note on arithmetic topology and dynamical systems,'' Contemp.~Math.~300, AMS (2002), pp.~99--114.

% --- p-adic physics ---
\bibitem{Volovich1987}
I.~V.~Volovich, ``Number theory as the ultimate physical theory,'' CERN-TH.4781/87 (1987); \textit{p-Adic Numbers Ultrametric Anal. Appl.} \textbf{2}, 77--87 (2010).

\bibitem{FreundWitten1987}
P.~G.~O.~Freund and E.~Witten, ``Adelic string amplitudes,'' \textit{Phys. Lett. B} \textbf{199}, 191--195 (1987).

\bibitem{FreundOlson1987}
P.~G.~O.~Freund and M.~Olson, ``Non-archimedean strings,'' \textit{Phys. Lett. B} \textbf{199}, 186--190 (1987).

\bibitem{Dragovich2009}
B.~Dragovich et al., ``On $p$-adic mathematical physics,'' \textit{p-Adic Numbers Ultrametric Anal. Appl.} \textbf{1}, 1--17 (2009). arXiv:0904.4205.

\bibitem{Manin1989}
Yu.~I.~Manin, ``Reflections on arithmetical physics,'' in \textit{Conformal Invariance and String Theory}, Academic Press (1989), pp.~293--303.

\bibitem{Manin1995}
Yu.~I.~Manin, ``Lectures on zeta functions and motives,'' \textit{Ast\'{e}risque} \textbf{228}, 121--163 (1995).

% --- Warp drive ---
\bibitem{Alcubierre1994}
M.~Alcubierre, ``The warp drive: hyper-fast travel within general relativity,'' \textit{Class. Quantum Grav.} \textbf{11}, L73--L77 (1994). arXiv:gr-qc/0009013.

\bibitem{Pfenning1997}
M.~J.~Pfenning and L.~H.~Ford, ``The unphysical nature of `warp drive','' \textit{Class. Quantum Grav.} \textbf{14}, 1743--1751 (1997). arXiv:gr-qc/9702026.

\bibitem{Lentz2021}
E.~W.~Lentz, ``Breaking the warp barrier: hyper-fast solitons in Einstein gravity,'' \textit{Class. Quantum Grav.} \textbf{38}, 075015 (2021). arXiv:2006.07125.

\bibitem{BobrickMartire2021}
A.~Bobrick and G.~Martire, ``Introducing physical warp drives,'' \textit{Class. Quantum Grav.} \textbf{38}, 105009 (2021). arXiv:2102.06824.

% --- Wormholes ---
\bibitem{VisserBook1995}
M.~Visser, \textit{Lorentzian Wormholes: From Einstein to Hawking}, AIP Press (1995).

% --- Topological quantum computation ---
\bibitem{Kitaev2003}
A.~Yu.~Kitaev, ``Fault-tolerant quantum computation by anyons,'' \textit{Ann. Phys.} \textbf{303}, 2--30 (2003). arXiv:quant-ph/9707021.

\bibitem{FreemanKitaevLarsenWang2003}
M.~H.~Freedman, A.~Kitaev, M.~J.~Larsen, and Z.~Wang, ``Topological quantum computation,'' \textit{Bull. Amer. Math. Soc.} \textbf{40}, 31--38 (2003). arXiv:quant-ph/0101025.

% --- Casimir effect ---
\bibitem{Casimir1948}
H.~B.~G.~Casimir, ``On the attraction between two perfectly conducting plates,'' \textit{Proc. K. Ned. Akad. Wet.} \textbf{51}, 793--795 (1948).

\bibitem{Lamoreaux1997}
S.~K.~Lamoreaux, ``Demonstration of the Casimir force in the 0.6 to 6 $\mu$m range,'' \textit{Phys. Rev. Lett.} \textbf{78}, 5--8 (1997).

\bibitem{Wilson2011}
C.~M.~Wilson et al., ``Observation of the dynamical Casimir effect in a superconducting circuit,'' \textit{Nature} \textbf{479}, 376--379 (2011).

% --- Additional references ---
\bibitem{AbramskyBrandenburger2011}
S.~Abramsky and A.~Brandenburger, ``The sheaf-theoretic structure of non-locality and contextuality,'' \textit{New J. Phys.} \textbf{13}, 113036 (2011). arXiv:1102.0264.

\bibitem{ConnesMarcolli2004}
A.~Connes and M.~Marcolli, ``Renormalization and motivic Galois theory,'' \textit{Int. Math. Res. Not.} \textbf{2004}(76), 4073--4091. arXiv:math/0409306.

\bibitem{Brown2012}
F.~Brown, ``Mixed Tate motives over $\mathbb{Z}$,'' \textit{Ann. of Math.} \textbf{175}, 949--976 (2012). arXiv:1102.1312.

\bibitem{Brown2017}
F.~Brown, ``Feynman amplitudes, coaction principle, and cosmic Galois group,'' \textit{Commun. Number Theory Phys.} \textbf{11}, 453--556 (2017). arXiv:1512.06409.

\bibitem{Marcolli2010}
M.~Marcolli, \textit{Feynman Motives}, World Scientific (2010).

\bibitem{BlochEsnaultKreimer2006}
S.~Bloch, H.~Esnault, and D.~Kreimer, ``On motives associated to graph polynomials,'' \textit{Commun. Math. Phys.} \textbf{267}, 181--225 (2006). arXiv:math/0510011.

\bibitem{Schreiber2013}
U.~Schreiber, ``Differential cohomology in a cohesive infinity-topos,'' arXiv:1310.7930 (2013).

\bibitem{Costello2011}
K.~Costello, \textit{Renormalization and Effective Field Theory}, AMS (2011).

\bibitem{ManinMarcolli2014}
Yu.~I.~Manin and M.~Marcolli, ``Big Bang, blowup, and modular curves,'' \textit{SIGMA} \textbf{10}, 073 (2014). arXiv:1402.2158.

\bibitem{Cartier2001}
P.~Cartier, ``A mad day's work: from Grothendieck to Connes and Kontsevich,'' \textit{Bull. Amer. Math. Soc.} \textbf{38}, 389--408 (2001).

\bibitem{Dragovich2017}
B.~Dragovich et al., ``$p$-adic mathematical physics: the first 30 years,'' \textit{p-Adic Numbers Ultrametric Anal. Appl.} \textbf{9}, 87--121 (2017). arXiv:1705.04758.

\bibitem{Dragovich1995}
B.~Dragovich, ``Adelic harmonic oscillator,'' \textit{Int. J. Mod. Phys. A} \textbf{10}, 2349--2365 (1995).

\bibitem{Nayak2008}
C.~Nayak et al., ``Non-Abelian anyons and topological quantum computation,'' \textit{Rev. Mod. Phys.} \textbf{80}, 1083--1159 (2008). arXiv:0707.1889.

\bibitem{Weinberg1989}
S.~Weinberg, ``The cosmological constant problem,'' \textit{Rev. Mod. Phys.} \textbf{61}, 1--23 (1989).

\bibitem{Mochizuki2020}
S.~Mochizuki, ``The mathematics of mutually alien copies: from Gaussian integrals to inter-universal Teichm\"{u}ller theory,'' RIMS preprint (2020).

\end{thebibliography}

\end{document}
